\documentclass[letterpaper,10pt]{article}

\usepackage[utf8]{inputenc}
\usepackage[spanish]{babel}
\usepackage{fontenc}
\usepackage[dvipdfmx]{graphicx}
\usepackage{bmpsize,wrapfig,xcolor}
\usepackage{fullpage}
\usepackage[hidelinks]{hyperref}

% Para evitar que se indente solo a cada rato
\setlength\parindent{0pt}

\begin{document}
	\begin{titlepage}

		\begin{wrapfigure}{R}{0.3\textwidth}
			\includegraphics[width=0.3\textwidth]{logoFCFM.png}
		\end{wrapfigure}

		\noindent \phantom - % "Hax" para que quede alineada la imagen con el texto

		Universidad de Chile

		Facultad de Ciencias Físicas y Matemáticas

		Depto. de Ciencias de la Computación

		CC5208 - Visualización de Información

		\vfill

		\begin{center}
			\begin{Huge}
				{\textbf{Tarea 3}}
			\end{Huge}

			\begin{large}
				Layout del proyecto
			\end{large}

		\end{center}

		\vfill

		\begin{flushright}
			\begin{tabular}{lll}
				Integrantes	&:	& Américo Ferrada\\
						&:	& Gabriel Sanhueza\\
				Profesor	&:	& Javier Bustos\\
				Ayudante	&:	& Diego Madariaga\\
			\end{tabular}
		\end{flushright}

	\end{titlepage}

	% % % % % % % % % % % % % % % % % % % % % % % % % % % % % % % % % % % % % % % % % % % % % % % % % % % % % % % % % % % % % % % % % % % % % % % % % % % % % % % % % % % % % % % % % %
	\newpage
	% % % % % % % % % % % % % % % % % % % % % % % % % % % % % % % % % % % % % % % % % % % % % % % % % % % % % % % % % % % % % % % % % % % % % % % % % % % % % % % % % % % % % % % % % %

	\tableofcontents

	% % % % % % % % % % % % % % % % % % % % % % % % % % % % % % % % % % % % % % % % % % % % % % % % % % % % % % % % % % % % % % % % % % % % % % % % % % % % % % % % % % % % % % % % % %
	\newpage
	% % % % % % % % % % % % % % % % % % % % % % % % % % % % % % % % % % % % % % % % % % % % % % % % % % % % % % % % % % % % % % % % % % % % % % % % % % % % % % % % % % % % % % % % % %

	\begin{large}
		{\textbf{BORRAR ESTO}}
	\end{large}

	En esta tarea debe enviar el diseño del layout para los datos de su proyecto. Es la última oportunidad que tiene para cambiar su dataset.

	Se espera que usted:

	\begin{enumerate}
		\item ** DONE: Haga un análisis de los tipos de datos que posee (ordinales, categóricos, cuantitativos, etc) y sus relaciones. **
		\item Plantee las preguntas que desea su visualización responda.
		\item Haga un mapeo de sus datos a los distintos gráficos que servirán para responder 2.
		\item Diseñe (usando lápiz/papel, paint, d3, tableau, lo que le sea más fácil) el layout de su visualización, textos de apoyo, y su paleta de colores,
		justificando sus decisiones de diseño desde el punto de vista cognitivo y con los fundamentos de información visual.
	\end{enumerate}

	Recuerde que hoy en clases se informó que sus datos y visualizaciones deben quedar en un repositorio de GitHub. Si usted no tiene uno propio debe avisarle al profesor para que le cree un repositorio.

	% % % % % % % % % % % % % % % % % % % % % % % % % % % % % % % % % % % % % % % % % % % % % % % % % % % % % % % % % % % % % % % % % % % % % % % % % % % % % % % % % % % % % % % % % %
	\newpage
	% % % % % % % % % % % % % % % % % % % % % % % % % % % % % % % % % % % % % % % % % % % % % % % % % % % % % % % % % % % % % % % % % % % % % % % % % % % % % % % % % % % % % % % % % %

	\section{Análisis}

	Existen 10 columnas de datos en total:

	\begin{itemize}
		\item Nombre: Aleatorio
		\item Id: Ordinal
		\item Tipo: Categórico
		\item Clase: Categórico
		\item Masa: Ordinal

		\item Caído/Encontrado: Categórico
		\item Año: Ordinal
		\item Latitud: Ordinal
		\item Longitud: Ordinal
		\item Geolocalización: Ordinal (Tupla)
	\end{itemize}

	Nuestros datos se relacionan de la siguiente manera:

	\begin{itemize}
		\item Tamaño y tipo del meteorito: Cada tipo de meteorito depende de su tamaño.
		\item Cantidad de meteoritos y año encontrado: Entre más reciente, mayor es la cantidad de meteoritos encontrados.
		\item Tamaño de los meteoritos y geolocalización: Se encuentra mayor cantidad de meteoritos pequeños en los polos.
		\item Nombre propio y tamaño del meteorito: Los meteoritos de tamaño no-despreciable tienen nombre propio (Ej.: Alessandria) v/s los más pequeños (Ej.: Yamato 983824)
	\end{itemize}

	% % % % % % % % % % % % % % % % % % % % % % % % % % % % % % % % % % % % % % % % % % % % % % % % % % % % % % % % % % % % % % % % % % % % % % % % % % % % % % % % % % % % % % % % % %
	\newpage
	% % % % % % % % % % % % % % % % % % % % % % % % % % % % % % % % % % % % % % % % % % % % % % % % % % % % % % % % % % % % % % % % % % % % % % % % % % % % % % % % % % % % % % % % % %

	\section{Planteamiento de preguntas}
	
	% FIXME Planteamiento de preguntas
	
	Preguntas que esperamos responder con nuestra Vizualizacion:
	
	\begin{itemize}
		\item ¿Donde se encuentran mas metioritos en funsión del peso?
		\item ¿Cual es el tipo de metiorito más frecuente?
		\item Lo anterior, pero respecto al peso.
		\item ¿Cual es el peso de metiorito mas frecuente?
		\item ¿Cual el el metiorito más pesado ncontrado?
		\item ¿Cual es el metiorito más liviano encontrado?
		\item ¿Que año se encontraron mas metioritos?
		\item ¿En que pais se encontraron mas metioritos?
		\item ¿Cual es el tipo de metiorito mas fracuente en funsión de la geolocalización?
	\end{itemize}

	% % % % % % % % % % % % % % % % % % % % % % % % % % % % % % % % % % % % % % % % % % % % % % % % % % % % % % % % % % % % % % % % % % % % % % % % % % % % % % % % % % % % % % % % % %
	\newpage
	% % % % % % % % % % % % % % % % % % % % % % % % % % % % % % % % % % % % % % % % % % % % % % % % % % % % % % % % % % % % % % % % % % % % % % % % % % % % % % % % % % % % % % % % % %

	\section{Mapeo datos a gráficos}
	% TODO Mapeo datos -> gráficos

	Como relacionar los datos para responder las oreguntas anteriores:
	
	\begin{itemize}
		\item Basta mapear los datos en un mapa del mundo respecto a su geolocalizacion, agregando un filtro en funsion del peso, secuenciando este dato en segmentos.
		\item En un graficode barra, por cada tipo contar las apariciones de cada tipo.
		\item Lo anterior, pero tomando solo un segmento de los datos en funcion del peso.
		\item En un grafico de barra ordenar los metioritos en funsion del peso y contar sus ocurrencias.
		\item En el grafico anterior se responden estas 3 preguntas.
		\item En un grafico de barra contar apariciones de metioritos en funsion del año.
		\item Dado un mapa del mundo, utilizar un gradiente de color en funsion del numero de metioritos caidos en un pais.
		\item Dado grupos de tipos similares, pintar un mapa en funsion de las ocurriencias del tipo.
	\end{itemize}

	% % % % % % % % % % % % % % % % % % % % % % % % % % % % % % % % % % % % % % % % % % % % % % % % % % % % % % % % % % % % % % % % % % % % % % % % % % % % % % % % % % % % % % % % % %
	\newpage
	% % % % % % % % % % % % % % % % % % % % % % % % % % % % % % % % % % % % % % % % % % % % % % % % % % % % % % % % % % % % % % % % % % % % % % % % % % % % % % % % % % % % % % % % % %

	\section{Diseño}
	% TODO Diseño

	TODO
	
	% 	Paleta de colores: http://paletton.com/#uid=50K0D0kS-C8tqT2FFTwQPquQols

\end{document}
